\documentclass[11pt]{article}

\usepackage[T1]{fontenc}
\usepackage{lmodern}
\usepackage{microtype}
\usepackage[margin=1in]{geometry}
\usepackage{amsmath,amssymb,amsthm,mathtools}
\usepackage{booktabs}
\usepackage{enumitem}
\usepackage{xcolor}
\usepackage{listings}
\usepackage{hyperref}
\usepackage[nameinlink,noabbrev]{cleveref}
\usepackage{fancyhdr}
\usepackage{lastpage}

\hypersetup{
  colorlinks=true,
  linkcolor=blue!55!black,
  citecolor=blue!55!black,
  urlcolor=blue!55!black,
  pdftitle={A variance-and-sieve route to Erdos Problem 971},
  pdfauthor={Kyungmin}
}

\pagestyle{fancy}
\fancyhf{}
\lhead{Candidate verification note for Erd\H{o}s Problem 971}
\rhead{Page \thepage\ of \pageref{LastPage}}
\setlength{\headheight}{14pt}

\newtheorem{theorem}{Theorem}[section]
\newtheorem{proposition}[theorem]{Proposition}
\newtheorem{lemma}[theorem]{Lemma}
\newtheorem{corollary}[theorem]{Corollary}
\theoremstyle{definition}
\newtheorem{definition}[theorem]{Definition}
\theoremstyle{remark}
\newtheorem{remark}[theorem]{Remark}
\newtheorem{checkpoint}[theorem]{Verification checkpoint}

\DeclareMathOperator{\li}{li}
\newcommand{\Pp}{\mathbb P}
\newcommand{\ph}{\varphi}
\newcommand{\RR}{\mathbb R}
\newcommand{\NN}{\mathbb N}
\newcommand{\ind}{\mathbf 1}
\newcommand{\Ss}{\mathfrak S}
\newcommand{\eps}{\varepsilon}
\newcommand{\bigO}{\mathrm O}

\definecolor{statusblue}{RGB}{235,244,252}
\definecolor{statusborder}{RGB}{48,97,145}
\definecolor{codebg}{RGB}{247,247,247}

\lstdefinelanguage{Lean}{
  morekeywords={theorem,lemma,def,namespace,end,noncomputable,section,open,import,structure,where,by,have,exact,apply,simpa,using,fun,if,then,else,let,in,with},
  sensitive=true,
  morecomment=[l]{--},
  morecomment=[s]{/-}{-/},
  morestring=[b]"
}
\lstset{
  language=Lean,
  basicstyle=\ttfamily\footnotesize,
  backgroundcolor=\color{codebg},
  frame=single,
  breaklines=true,
  columns=fullflexible,
  keepspaces=true,
  showstringspaces=false,
  xleftmargin=0.5em,
  xrightmargin=0.5em
}

\title{\bfseries A Variance-and-Sieve Route to Erd\H{o}s Problem 971\\[0.5em]
\large Candidate proof and Lean-checked finite reduction}
\author{Kyungmin}
\date{25 July 2026}

\begin{document}
\maketitle

\begin{center}
\fcolorbox{statusborder}{statusblue}{%
\begin{minipage}{0.91\textwidth}
\textbf{Submission status.}
This is a candidate unconditional argument posted for independent verification, not a claim that the problem has already been accepted as solved.
The accompanying Lean 4 file kernel-checks the finite combinatorial reduction and the bridge from prime occupancies to least congruent primes, conditional on explicit analytic estimates.
It does \emph{not} formalize the Friedlander--Goldston variance theorem, the uniform three-tuple upper sieve, or the asymptotic estimates used to instantiate those hypotheses.
A successful Lean compilation therefore verifies the logical core, but not yet the full analytic-number-theory proof.
\end{minipage}}
\end{center}

\begin{abstract}
Let $p(a,q)$ denote the least prime congruent to $a \pmod q$.
Erd\H{o}s Problem 971 asks whether there is an absolute $c>0$ such that, for every sufficiently large $q$,
\[
  p(a,q)>(1+c)\ph(q)\log q
\]
for a positive proportion of the reduced residue classes $a\pmod q$.
This note presents a candidate argument based on two published sieve-theoretic inputs.
At the critical cutoff $X\sim \ph(q)\log q$, a pointwise lower bound for Hooley's variance yields a lower bound $C_2\gg\ph(q)$ for the second factorial moment of the numbers of primes in the reduced residue classes.
A uniform upper-bound sieve for three linear forms, combined with an elementary average of its singular series, yields $C_3\ll\ph(q)$ for the third factorial moment.
These two estimates imply that a positive proportion of classes contain at least two primes below $X$.
Since the total number of prime incidences below $X$ is only $\ph(q)+o(\ph(q))$, a positive proportion of classes are empty.
The empty-class estimate survives enlargement of the cutoff to $(1+c)\ph(q)\log q$ for sufficiently small fixed $c>0$.
The accompanying Lean development checks the entire finite implication from the moment hypotheses to the least-prime conclusion.
The main purpose of this submission is to expose the exact analytic and formal verification points to expert scrutiny.
\end{abstract}

\tableofcontents

\section{Problem and claimed conclusion}

For $(a,q)=1$, define
\[
  p(a,q):=\min\{p\in\Pp:p\equiv a\pmod q\}.
\]
Dirichlet's theorem guarantees that the set is nonempty.
The problem asks whether there are absolute constants $c,C>0$ such that, for all sufficiently large $q$,
\begin{equation}\label{eq:target}
  \#\left\{a\bmod q:(a,q)=1,\ p(a,q)>(1+c)\ph(q)\log q\right\}
  \ge C\ph(q).
\end{equation}
The Erd\H{o}s Problems database currently lists this as open and records that Erd\H{o}s proved the assertion only for an infinite sequence of moduli \cite{ErdosProblems971,Erdos1950}.
Leung obtained the expected exponential law for the least prime under a uniform Hardy--Littlewood prime $k$-tuple hypothesis \cite{Leung2024}.

The argument below proposes an unconditional proof of \eqref{eq:target} from existing pointwise variance and upper-bound sieve estimates.
Because the conclusion would close a problem still publicly listed as open, every uniformity statement is isolated explicitly.

\section{Critical cutoff and factorial moments}

Write
\[
  \Phi:=\ph(q),\qquad
  R:=\left\lfloor\frac{\Phi\log q}{q}\right\rfloor,
  \qquad X:=qR.
\]
The standard lower bound $\ph(q)/q\gg 1/\log\log q$ implies $R\to\infty$.
Moreover,
\begin{equation}\label{eq:cutoff-asymp}
  X=(1+o(1))\Phi\log q,
  \qquad \log X=(1+o(1))\log q,
  \qquad \frac{X}{\log X}=(1+o(1))\Phi.
\end{equation}
Since $R\le \log q$ and $X\ge q$ for large $q$,
\begin{equation}\label{eq:FG-range}
  q=\frac XR\ge \frac{X}{\log X}.
\end{equation}
Thus $(X,q)$ lies in the large-modulus range required by the Friedlander--Goldston lower bound with $A=1$.

For each reduced residue class $a\pmod q$, set
\[
  N_a=N_a(X):=\#\{p\le X:p\equiv a\pmod q\}.
\]
Define the factorial moments
\begin{equation}\label{eq:moments}
  C_2:=\sum_{a\bmod q}^{*}\binom{N_a}{2},
  \qquad
  C_3:=\sum_{a\bmod q}^{*}\binom{N_a}{3}.
\end{equation}
The proposed proof reduces the analytic work to
\begin{equation}\label{eq:key-moments}
  C_2\gg\Phi,
  \qquad
  C_3\ll\Phi.
\end{equation}

\section{The second moment from pointwise variance}

For a reduced class $a\pmod q$, put
\[
  \psi_a:=\sum_{\substack{n\le X\\ n\equiv a\ (q)}}\Lambda(n),
  \qquad
  S:=\sum_{\substack{n\le X\\(n,q)=1}}\Lambda(n).
\]
Use the naturally centered variance
\begin{equation}\label{eq:variance}
  V_S(X;q):=\sum_{a\bmod q}^{*}
  \left(\psi_a-\frac{S}{\Phi}\right)^2.
\end{equation}
Friedlander and Goldston proved a pointwise lower bound of the conjectured order in the range $X/(\log X)^A\le q\le X$; Fiorilli records it in the form
\begin{equation}\label{eq:FG}
  V(X;q)\ge \left(\frac12-o(1)\right)X\log q
\end{equation}
for every fixed $A>0$, uniformly in that range \cite{FriedlanderGoldston1996,Fiorilli2015}.
If the cited formulation is centered at $X/\Phi$ rather than $S/\Phi$, the identity
\[
  \sum_a^*\left(\psi_a-\frac{X}{\Phi}\right)^2
  =V_S(X;q)+\frac{(S-X)^2}{\Phi}
\]
and $S=X+o(X)$ transfer \eqref{eq:FG} to \eqref{eq:variance}, because $(S-X)^2/\Phi=o(X\log q)$ at the present cutoff.

Define
\[
  D_\Lambda:=\sum_{\substack{n\le X\\(n,q)=1}}\Lambda(n)^2
\]
and
\[
  P_\Lambda:=
  \sum_{\substack{m<n\le X\\m\equiv n\ (q)\\(mn,q)=1}}
  \Lambda(m)\Lambda(n).
\]
Expanding the square gives the exact identity
\begin{equation}\label{eq:variance-expansion}
  V_S(X;q)=D_\Lambda+2P_\Lambda-\frac{S^2}{\Phi}.
\end{equation}
The prime number theorem and partial summation give
\[
  S=X+o(X),\qquad D_\Lambda=X\log X+O(X).
\]
Deleting prime powers whose underlying prime divides $q$ costs only $O(\log X\log q)$.
From the definition of $X$,
\[
  \frac{X}{\Phi}=\log q+O\!\left(\frac q\Phi\right)
  =\log q+O(\log\log q),
\]
and $\log X=\log q+O(\log\log q)$. Hence
\begin{equation}\label{eq:diagonal-cancel}
  D_\Lambda-\frac{S^2}{\Phi}=o(X\log q).
\end{equation}
Combining \eqref{eq:FG}, \eqref{eq:variance-expansion}, and \eqref{eq:diagonal-cancel} yields
\begin{equation}\label{eq:weighted-pairs}
  P_\Lambda\ge \left(\frac14-o(1)\right)X\log q.
\end{equation}

There are $O(X^{1/2})$ proper prime powers below $X$.
For each such number there are at most $R+1=O(\log X)$ numbers below $X$ in the same residue class, and each pair has weight at most $(\log X)^2$.
Thus all pairs involving a proper prime power contribute
\[
  O\!\left(X^{1/2}R(\log X)^2\right)=o(X\log q).
\]
The remaining contribution in \eqref{eq:weighted-pairs} comes from pairs of distinct primes.
Since every prime-pair weight is at most $(\log X)^2$,
\begin{equation}\label{eq:C2}
  C_2\gg \frac{X\log q}{(\log X)^2}\gg\Phi.
\end{equation}

\begin{checkpoint}[Pointwise variance input]\label{check:variance}
The crucial issue is that \eqref{eq:FG} must be a lower bound for each individual modulus $q$, uniformly for $X/(\log X)^A\le q\le X$, rather than an estimate only after averaging over $q$.
Fiorilli's introduction explicitly distinguishes the average results from the Friedlander--Goldston pointwise lower bound and states the latter with constant $1/2$ \cite{Fiorilli2015}.
Reviewers should nevertheless compare the exact normalization and uniformity with the original theorem in \cite{FriedlanderGoldston1996}.
\end{checkpoint}

\section{The third moment from a uniform three-tuple upper sieve}

Each unordered triple $p_0<p_1<p_2\le X$ in one residue class has a unique representation
\[
  p_1=p_0+rq,\qquad p_2=p_0+sq,
  \qquad 1\le r<s<R.
\]
Therefore
\begin{equation}\label{eq:C3-shifts}
  C_3=\sum_{1\le r<s<R}T_q(r,s),
\end{equation}
where
\[
  T_q(r,s):=
  \#\{p\le X-sq:p,\ p+rq,\ p+sq\text{ are prime}\}.
\]
We may enlarge the range to $p\le X$.

For a prime $\ell$, let $\nu_\ell(r,s;q)$ be the number of distinct residues among
\[
  0,\ -rq,\ -sq\pmod\ell.
\]
Define
\begin{equation}\label{eq:singular-series}
  \Ss_q(r,s):=
  \prod_{\ell}
  \left(1-\frac{\nu_\ell(r,s;q)}{\ell}\right)
  \left(1-\frac1\ell\right)^{-3}.
\end{equation}
If the triple is admissible, a uniform upper-bound sieve gives
\begin{equation}\label{eq:tuple-upper}
  T_q(r,s)\ll \Ss_q(r,s)\frac{X}{(\log X)^3},
\end{equation}
with an absolute implied constant.
For direct reproducibility, Theorem 1 of Dubbe gives an explicit $K=3$ version for primes $p\le X$ such that $a_i p+b_i$ are prime; taking $a_1=a_2=1$, $b_1=rq$, and $b_2=sq$ gives \eqref{eq:tuple-upper} \cite{Dubbe2024}.
The nondegeneracy determinant is $r s(s-r)q^3\ne0$.

If the triple is not admissible, some prime $\ell\le3$ divides one of the three values for every $p$.
If all three values are prime, one of them must equal $\ell$.
Consequently
\begin{equation}\label{eq:nonadmissible}
  T_q(r,s)\le3
\end{equation}
for every nonadmissible pair $(r,s)$.

\subsection{Uniform singular-series average}

If $\ell\mid q$, then $\nu_\ell=1$, and the local factor in \eqref{eq:singular-series} is
\[
  \left(1-\frac1\ell\right)^{-2}.
\]
Thus the full contribution of primes dividing $q$ is exactly
\begin{equation}\label{eq:q-factor}
  \prod_{\ell\mid q}\left(1-\frac1\ell\right)^{-2}
  =\left(\frac q\Phi\right)^2.
\end{equation}

Suppose now that $\ell\nmid q$.
If $\ell\nmid rs(s-r)$, then $\nu_\ell=3$ and, for $\ell\ge5$,
\[
  \left(1-\frac3\ell\right)\left(1-\frac1\ell\right)^{-3}<1.
\]
Hence generic primes can be discarded in an upper bound.
If $\ell\mid rs(s-r)$, then $\nu_\ell\in\{1,2\}$ and the local factor is at most $1+K/\ell$ for an absolute $K$; the finitely many local possibilities at $2$ and $3$ are absorbed into the implied constant.
Set
\[
  F(n):=\prod_{\ell\mid n}\left(1+\frac K\ell\right).
\]
Then
\begin{equation}\label{eq:series-majorant}
  \Ss_q(r,s)\ll
  \left(\frac q\Phi\right)^2F(r)F(s)F(s-r).
\end{equation}

For a sufficiently large absolute $K_1$,
\[
  F(n)^3\le
  \prod_{\ell\mid n}\left(1+\frac{K_1}{\ell}\right)
  =\sum_{d\mid n}\frac{\mu^2(d)K_1^{\omega(d)}}{d}.
\]
Therefore
\begin{equation}\label{eq:Fmean}
  \sum_{n\le R}F(n)^3
  \le R\sum_{d\ge1}\frac{\mu^2(d)K_1^{\omega(d)}}{d^2}
  \ll R.
\end{equation}
Applying H\"older's inequality over the pairs $1\le r<s<R$ gives
\begin{equation}\label{eq:holder}
  \sum_{1\le r<s<R}F(r)F(s)F(s-r)\ll R^2.
\end{equation}
Consequently
\begin{equation}\label{eq:series-average}
  \sum_{1\le r<s<R}\Ss_q(r,s)
  \ll R^2\left(\frac q\Phi\right)^2.
\end{equation}

Combining \eqref{eq:C3-shifts}, \eqref{eq:tuple-upper}, \eqref{eq:nonadmissible}, and \eqref{eq:series-average},
\begin{align*}
  C_3
  &\ll \frac{X}{(\log X)^3}
       R^2\left(\frac q\Phi\right)^2+R^2\\
  &\ll \Phi.
\end{align*}
Indeed $R\le \Phi\log q/q$, $X\asymp\Phi\log q$, and $\log X\asymp\log q$.
This proves the second estimate in \eqref{eq:key-moments}.

\begin{checkpoint}[Three-tuple sieve and local factors]\label{check:triple}
The two points most deserving independent checking are:
\begin{enumerate}[label=(\roman*),leftmargin=2.2em]
  \item that the upper bound \eqref{eq:tuple-upper} is uniform when the shifts $rq,sq$ vary with $q$ and can be as large as $X$; and
  \item that the factor at every prime dividing $q$ is exactly $(1-1/\ell)^{-2}$, so that the global factor $(q/\ph(q))^2$ is cancelled by the number $R^2$ of shift pairs.
\end{enumerate}
The explicit formulation in \cite[Theorem 1]{Dubbe2024} appears to supply the required uniformity.
\end{checkpoint}

\section{Finite extraction of many repeated and empty classes}

Choose absolute constants $\alpha,\beta>0$ such that, for all sufficiently large $q$,
\[
  C_2\ge\alpha\Phi,
  \qquad C_3\le\beta\Phi.
\]
Let
\[
  D_2:=\#\{a\bmod q:(a,q)=1,\ N_a\ge2\}.
\]
Fix an integer $T\ge3$ so large that
\begin{equation}\label{eq:Tchoice}
  \frac{3\beta}{T-1}\le\frac\alpha2.
\end{equation}
For $k>T$,
\[
  \binom{k}{2}=\frac{3}{k-2}\binom{k}{3}
  \le\frac{3}{T-1}\binom{k}{3}.
\]
Hence the contribution to $C_2$ from classes with $N_a>T$ is at most $\alpha\Phi/2$.
It follows that
\[
  \sum_{2\le N_a\le T}\binom{N_a}{2}\ge\frac\alpha2\Phi.
\]
Each summand is at most $\binom T2$, and therefore
\begin{equation}\label{eq:D2}
  D_2\ge \delta\Phi,
  \qquad
  \delta:=\frac{\alpha}{2\binom T2}>0.
\end{equation}

Let
\[
  M:=\sum_{a\bmod q}^{*}N_a.
\]
For large $q$, all prime divisors of $q$ lie below $X$, so
\[
  M=\pi(X)-\omega(q)=\Phi+o(\Phi)
\]
by \eqref{eq:cutoff-asymp} and the prime number theorem.
Let $O(X)$ denote the number of occupied reduced residue classes and $U(X)=\Phi-O(X)$ the number of empty ones.
Every occupied class consumes one prime incidence, and every class counted by $D_2$ consumes at least one additional incidence. Thus
\[
  M\ge O(X)+D_2.
\]
Using \eqref{eq:D2},
\begin{equation}\label{eq:empty-base}
  U(X)\ge \Phi-M+D_2\ge(\delta-o(1))\Phi.
\end{equation}

This is the finite combinatorial step checked by the accompanying Lean development.

\section{Enlarging the cutoff and completing the candidate proof}

Let
\[
  Y:=(1+c)\Phi\log q
\]
for a fixed $c>0$ to be chosen.
Since $X=(1+o(1))\Phi\log q$, the prime number theorem gives
\begin{equation}\label{eq:new-primes}
  \pi(Y)-\pi(X)=c\Phi+o(\Phi).
\end{equation}
Each newly admitted prime can fill at most one class that was empty at $X$.
Therefore
\[
  U(Y)\ge U(X)-\bigl(\pi(Y)-\pi(X)\bigr)
  \ge(\delta-c-o(1))\Phi.
\]
Choose, for example, $c=\delta/4$.
Then for all sufficiently large $q$,
\[
  U(Y)\ge\frac\delta2\Phi.
\]
An empty class up to $\lfloor Y\rfloor$ has least congruent prime greater than $Y$.
This gives \eqref{eq:target} with $C=\delta/2$.

\begin{theorem}[Candidate conclusion]\label{thm:candidate}
If the pointwise variance lower bound used in \cref{check:variance} and the uniform three-tuple upper sieve used in \cref{check:triple} apply with the stated uniformity, then there exist absolute constants $c,C>0$ such that \eqref{eq:target} holds for every sufficiently large $q$.
\end{theorem}

\section{Lean formalization and its exact scope}

The file \texttt{lean/Erdos971Forum.lean} is a standalone Lean 4 development importing \texttt{Mathlib}.
It defines finite occupancies, second and third factorial moments, empty and repeated class counts, prime occupancies modulo $d$, and the least prime in a congruence class.
Its strongest arithmetic specialization is:

\begin{lstlisting}
theorem prime_moment_method_to_least_prime_classes
    {d X0 X1 T : Nat} (hd : d != 0) (hT : 3 <= T) (hcut : X0 <= X1)
    (pairConst tripleConst doubleConst massError extensionConst : Real)
    (hPair : ...)
    (hTriple : ...)
    (hGap : ...)
    (hTotal : ...)
    (hIncrement : ...) :
    (doubleConst - massError - extensionConst) * (d.totient : Real) <=
      ((emptyLeastPrimeClasses d X1).card : Real)
\end{lstlisting}

The omitted hypotheses are written in full in the source file. They are exactly:
\begin{enumerate}[leftmargin=2.2em]
  \item a lower bound for the second factorial moment;
  \item an upper bound for the third factorial moment;
  \item a scalar truncation inequality selecting a fixed occupancy threshold;
  \item an upper bound for total prime incidence at the base cutoff; and
  \item an upper bound for incidences added when the cutoff is enlarged.
\end{enumerate}
The proof then returns a lower bound for the cardinality of
\[
  \{a<d:(a,d)=1,\ X_1<p(a,d)\}.
\]
The bridge from zero occupancy to the least-prime inequality uses mathlib's formalization of Dirichlet's theorem on primes in arithmetic progressions.

\subsection{Reproducible checks}

The package is pinned to Lean 4.27.0 and mathlib 4.27.0.
From the \texttt{lean/} directory, run
\begin{lstlisting}[language=bash]
./check_source.sh
\end{lstlisting}
or equivalently
\begin{lstlisting}[language=bash]
lake update
lake build
lake env lean AnalyticTargets.lean
lake env lean AuditAxioms.lean
\end{lstlisting}
The script rejects textual occurrences of \texttt{sorry}, \texttt{admit}, and \texttt{axiom} in the compiled core, and fails if \texttt{sorryAx} appears in the printed theorem dependencies.
Foundational dependencies such as \texttt{propext}, \texttt{Quot.sound}, or \texttt{Classical.choice} can be normal in mathlib developments.

\begin{center}
\fcolorbox{statusborder}{statusblue}{%
\begin{minipage}{0.91\textwidth}
\textbf{What a successful Lean run establishes.}
It establishes that the finite moment hypotheses imply the stated lower bound for classes with large least congruent prime, and that this implication contains no admitted Lean proof gaps.
It does \emph{not} establish the analytic moment hypotheses themselves.
Those remain the principal target for a complete end-to-end formal proof.
\end{minipage}}
\end{center}

\section{Requested expert checks}

A useful review can focus on the following finite list rather than rechecking every elementary step.

\begin{enumerate}[leftmargin=2.4em]
  \item \textbf{Variance theorem.} Verify the definition, centering, and uniformity of the Friedlander--Goldston bound for $q\ge X/\log X$.
  \item \textbf{Diagonal cancellation.} Confirm that \eqref{eq:diagonal-cancel} is uniform for arbitrary, possibly very smooth, $q$.
  \item \textbf{Prime powers.} Confirm the $o(X\log q)$ removal estimate in the weighted pair sum.
  \item \textbf{Upper sieve.} Confirm that Dubbe's theorem, or a standard Selberg upper-bound sieve theorem, applies uniformly to $p,p+rq,p+sq$ for $r,s<R$.
  \item \textbf{Nonadmissible patterns.} Confirm the bound $T_q(r,s)\le3$.
  \item \textbf{Singular series.} Check the local-factor majorant \eqref{eq:series-majorant} and the H\"older average \eqref{eq:series-average}, especially when $q$ is highly composite.
  \item \textbf{Cutoff bookkeeping.} Check floors, strict inequalities, and the passage from $X$ to the real cutoff $Y=(1+c)\Phi\log q$.
  \item \textbf{Lean scope.} Rebuild the package and inspect \texttt{AuditAxioms.lean}; do not interpret compilation of the conditional core as a formal proof of the analytic inputs.
\end{enumerate}

Any failure in items 1 or 4--6 would invalidate the proposed closing argument. Conversely, if all eight checks pass, the argument appears to give an unconditional affirmative answer to Erd\H{o}s Problem 971.

\section*{AI assistance disclosure}
\addcontentsline{toc}{section}{AI assistance disclosure}

The mathematical investigation, development of the candidate proof strategy,
and preparation of this verification package were conducted with substantial
assistance from OpenAI's GPT-5.6 Pro. The named human author is responsible for
the decision to submit this draft and for its claims. The use of an AI system
does not replace independent checking of the cited analytic inputs or expert
review of the proposed argument.

\section*{Acknowledgment of uncertainty}
\addcontentsline{toc}{section}{Acknowledgment of uncertainty}

This note was prepared as a verification-oriented research draft.
The fact that the public problem remains open is substantial contrary evidence against treating the argument as settled without expert review.
The Lean artifact narrows the possible failure to the analytic interface; it does not remove that possibility.

\begin{thebibliography}{99}

\bibitem{Dubbe2024}
T. Dubbe,
\emph{Explicit bounds for prime $k$-tuplets},
arXiv:2409.04261, 2024.
\url{https://arxiv.org/abs/2409.04261}

\bibitem{Erdos1950}
P. Erd\H{o}s,
\emph{On some applications of Brun's method},
Acta Sci. Math. (Szeged) \textbf{13} (1950), no. 1, 57--63.

\bibitem{ErdosProblems971}
T. F. Bloom,
\emph{Erd\H{o}s Problem 971},
Erd\H{o}s Problems database, accessed 25 July 2026.
\url{https://www.erdosproblems.com/971}

\bibitem{Fiorilli2015}
D. Fiorilli,
\emph{The distribution of the variance of primes in arithmetic progressions},
Int. Math. Res. Not. IMRN 2015, no. 12, 4421--4448;
arXiv:1301.5663.
\url{https://arxiv.org/abs/1301.5663}

\bibitem{FriedlanderGoldston1996}
J. B. Friedlander and D. A. Goldston,
\emph{Variance of distribution of primes in residue classes},
Q. J. Math. \textbf{47} (1996), 313--336.
\url{https://doi.org/10.1093/qmath/47.3.313}

\bibitem{FormalConjectures}
The Formal Conjectures Authors,
\emph{Formal Conjectures repository}, 2025--2026.
\url{https://github.com/google-deepmind/formal-conjectures}

\bibitem{Leung2024}
S.-K. Leung,
\emph{Moments of primes in progressions to a large modulus},
arXiv:2402.07941, 2024.
\url{https://arxiv.org/abs/2402.07941}

\end{thebibliography}

\end{document}
